\documentclass[11pt]{article}
\usepackage[a4paper,margin=24mm]{geometry}
\usepackage{amsmath,amssymb,amsthm,mathtools,mathrsfs,graphicx,booktabs}
\usepackage[colorlinks=true,linkcolor=blue,urlcolor=blue,citecolor=blue]{hyperref}
\usepackage{microtype}
\newtheorem{theorem}{Theorem}[section]
\newtheorem{lemma}[theorem]{Lemma}
\newtheorem{proposition}[theorem]{Proposition}
\DeclareMathOperator{\Res}{Res}
\DeclareMathOperator{\rem}{rem}
\DeclareMathOperator{\Tr}{Tr}
\DeclareMathOperator{\diag}{diag}
\newcommand{\C}{\mathbb C}
\newcommand{\R}{\mathbb R}
\newcommand{\B}{\mathscr B}
\newcommand{\MM}{\mathcal M}
\newcommand{\PP}{\operatorname{Prim}_h}
\newcommand{\HS}{\mathsf H_c}
\newcommand{\dd}{\,d}
\title{Constructive two-prime packets:\\exact theta lifts and bounded moment dynamics}
\author{Programme calculation; private working note}
\date{23 September 2026}
\begin{document}
\maketitle
\begin{abstract}
The original prime actions for 2 and 3 recover a finite zero packet with
every multiplicity. We turn that generation statement into a finite
construction: a bounded search for a real generator, an invertible
full-jet matrix, an explicit polynomial inverse, and prime words of
length at most one less than the packet's total multiplicity.
The same words act by actual dilations of the original theta source.
Their exact cochain discrepancy is calculated, rather than discarded.
Their inherited $L^2(dx)$ Gram matrices come from a compactly supported
arithmetic measure, whose recurrence coefficients have a bound
independent of the cutoff. The associated source/relation determinant
ratio satisfies an explicit bounded Toda curvature identity.
An original-zeta example is evaluated numerically. No off-critical zero
is invented or found, no positivity of the full Weil form is assumed,
and no existing Euler-polynomial metric is replaced by the new word metric.
\end{abstract}

\section{Original objects, source identity, and the calculation being extended}
The immediate programme inputs are TP1--8 (the two-prime torus receiver)
and PL1--27 (primitive full-jet representatives and their exact theta
relations) in the current, unsealed source of the task
\emph{Split Zero to GCT bridge}. Their private proof files and byte hashes
are identified in \texttt{SOURCE\_READING.json}; they have not been uploaded
by this calculation. We prove the ingredients needed here below.
The human Mellin-quotient context is Ralf Meyer's original
\href{https://arxiv.org/abs/math/0412277}{spectral interpretation},
Theorem \texttt{the:Lap\_range} and Corollary
\texttt{cor:Dminus\_eigenvalues} \cite{Meyer}. The interpolation step is
classical Chinese remainder/Hermite interpolation; its human algebraic
source is \href{https://stacks.math.columbia.edu/tag/00DT}{Stacks tag 00DT}
\cite{Stacks}. The specialised coefficients, lifts, and bounds are
derived explicitly here, not attributed to those sources as new results.

Keep
\begin{align}
 V&=\{\phi\in\mathcal S(\R):\phi\text{ even},\ \phi(0)=0,
                         \ \int_\R\phi(x)\dd x=0\},\nonumber\\
 \B&=\{F\in C^\infty(\R_{>0}):
       \sup_{x>0}|x^bD^jF(x)|<\infty\quad(b\in\mathbb Z,j\ge0)\},\nonumber\\
 D&=-x\partial_x,\qquad\Theta\phi(x)=2\sum_{n\ge1}\phi(nx),
 \qquad Q=\B/\Theta V.\label{eq:spaces}
\end{align}
The observation is the actual measure $dx$ on $\R_{>0}$, not $dx/x$.
The Mellin transform itself uses $dx/x$:
\begin{equation}
 \MM F(s)=\int_0^\infty F(x)x^s\frac{dx}{x},\qquad
 \|F\|_{L^2(dx)}^2=\frac1{2\pi}\int_\R
                         |\MM F(1/2+it)|^2\dd t.\label{eq:plancherel}
\end{equation}
Indeed $x=e^y$ gives the ordinary Fourier transform of
$e^{y/2}F(e^y)$, with kernel $e^{ity}$; Fourier--Plancherel proves
the equality with the factor $1/(2\pi)$ retained.

Write
\begin{equation}
 a(s)=s(s-1),\quad
 g(s)=2\xi(s)=a(s)\pi^{-s/2}\Gamma(s/2)\zeta(s),\quad
 \phi_*(x)=(4\pi^2x^4-6\pi x^2)e^{-\pi x^2}.\label{eq:seed}
\end{equation}
Here $g$ is the entire completion; Meyer's notation for a completed
function must not be substituted for this definition. Direct Gaussian
integration gives $\phi_*\in V$ and
$M_+\phi_*=a(s)\pi^{-s/2}\Gamma(s/2)/2$.
Absolute summation for $\Re s>1$ gives
$\MM\Theta\phi=2\zeta M_+\phi$, hence
$\MM\Theta\phi_*=g$; continuation gives these identities wherever used.
Poisson summation and the two zero moments show $\Theta V\subset\B$.
The Fourier convention in this step is $e^{-2\pi ix\xi}$.
Also $\MM DF=s\MM F$ by integration by parts.

Let $Z$ be a nonempty finite set of actual nontrivial zeros, stable under
$\iota\rho=1-\bar\rho$, with their full orders $m_\rho$.
The leading coefficient below is not changed:
\begin{equation}
 h(s)=c_h\prod_{\rho\in Z}(s-\rho)^{m_\rho},\quad c_h\ne0,
 \quad d=\sum_\rho m_\rho,\quad U_h=g/h,\quad
 A_h=\C[s]/(h)=\prod_\rho\C[t_\rho]/(t_\rho^{m_\rho}).\label{eq:packet}
\end{equation}
The coordinate $t_\rho=s-\rho$ and the conjugation
$f^\dagger(s)=\overline{f(1-\bar s)}$ remain fixed.
Both $[a]$ and $[U_h]$ are units in this full, possibly nonreduced algebra.
For $u\in A_h$ put
\begin{equation}
 R_u=\rem_h([aU_h]^{-1}u),\qquad
 \PP(u)=a(D)S_h\Theta(R_u(D)\phi_*),\qquad
 \MM\PP(u)=aU_hR_u.\label{eq:primitive}
\end{equation}
Here the inverse Euler operators are actual integrals:
\begin{equation}
 S_\rho H(x)=x^{-\rho}\int_x^\infty H(t)t^{\rho-1}\dd t,
 \qquad S_h=c_h^{-1}\prod_\rho S_\rho^{m_\rho}.\label{eq:inverse}
\end{equation}
They are applied only after the corresponding full zero orders are
present in the Mellin transform. Differentiation gives
$(D-\rho)S_\rho H=H$. At zero, the vanishing Mellin moment changes
the tail to minus the integral from zero to $x$. Rapid decay of $H$
then gives rapid decay at both ends for $S_\rho H$ and all Euler
derivatives. For example, bounds $|H(e^y)|\le C_Ne^{-Ny}$ on $y\ge0$
and $|H(e^y)|\le C'_Ne^{Ny}$ on $y\le0$ give respectively denominators
$N-\Re\rho$ and $N+\Re\rho$ when these are positive.
Integration by parts proves its Mellin quotient. Repeating removes
each required order. Injectivity of every $D-\rho$ on $\B$ (its
nonzero homogeneous solution fails rapid decay at an end) proves
independence of the order of these inverse operations.

Thus \eqref{eq:primitive} belongs to $\B$ and has endpoint Mellin values
zero. Its full jets at $Z$ equal $u$, and at every other zero of $g$
they vanish to the full order. Its class in $Q$ is injective in $u$:
a theta relation has zero selected jets. These facts establish the
literal map that will receive the prime words, without assuming a
positive form on the quotient.

\section{A finite integer-selected real generator with every original jet}\label{sec:constructive-generator}

\subsection{Original objects and the integer selector}

Let $Z=\{\rho_1,\ldots,\rho_r\}$ be a nonempty finite set of actual
zeros of $g=2\xi$, invariant under $\iota\rho=1-\bar\rho$.
Retain each full multiplicity $m_\rho\geq1$, so that
$m_{\iota\rho}=m_\rho$, and put $d=\sum_{\rho\in Z}m_\rho$.
Keep the original local coordinates and reflection:
\begin{equation}
 A_Z=\prod_{\rho\in Z}\mathbb{C}[t_\rho]/(t_\rho^{m_\rho}),\qquad
 t_\rho=s-\rho,\qquad
 (f^\dagger)_{\rho,j}=(-1)^j\overline{f_{\iota\rho,j}},\qquad
 A_{Z,\mathbb{R}}=A_Z^{\dagger}.
 \tag{CG1}\label{cg:original}
\end{equation}
The relation $f=f^\dagger$ is therefore
$f_{\iota\rho,j}=(-1)^j\overline{f_{\rho,j}}$.
No original coordinate or multiplicity is changed below.
For $p=2,3$ write
\begin{equation}
 \lambda_p(s)=p^{s-1/2},\qquad
 a_p(s)=\frac{\lambda_p(s)+\lambda_p(s)^{-1}}2,\qquad
 b_p(s)=\frac{\lambda_p(s)-\lambda_p(s)^{-1}}{2i}.
 \tag{CG2}\label{cg:prime}
\end{equation}
Their images in $A_Z$ are their full Taylor expansions through
$t_\rho^{m_\rho-1}$. The original prime action is
$\Pi_p=M_{p^s}=p^{1/2}M_{\lambda_p}$, not $M_{\lambda_p}$.
The identities $\lambda_p^\dagger=\lambda_p^{-1}$,
$a_p^\dagger=a_p$, $b_p^\dagger=b_p$, and $a_p^2+b_p^2=1$
hold with these exact expressions.

For an integer $c$ set
\begin{equation}
 H_c=a_2+c b_2+c^2a_3+c^3b_3,\qquad
 z_\rho=H_c(\rho),\qquad d_\rho=H_c'(\rho).
 \tag{CG3}\label{cg:generator}
\end{equation}
Here $H_c$ also denotes the entire function in the original variable
$s$; its class is taken only when it is used in $A_Z$.

\begin{theorem}[A bounded finite integer choice]\label{cg:selector-thm}
Put $r_2=\#\{\rho\in Z:m_\rho>1\}$ and
\begin{equation}
 B=3\left(\binom r2+r_2\right).
 \tag{CG4}\label{cg:bound}
\end{equation}
At least one $c\in\{0,1,\ldots,B\}$ has distinct values $z_\rho$
at all distinct roots and $d_\rho\ne0$ at every root of multiplicity greater than one. Every such
$c$ gives $\mathbb{R}[H_c]=A_{Z,\mathbb{R}}$ and $\mathbb{C}[H_c]=A_Z$.
No critical-line or simple-zero hypothesis is needed.
\end{theorem}

\begin{proof}
If both $\lambda_2(\rho)=\lambda_2(\sigma)$ and
$\lambda_3(\rho)=\lambda_3(\sigma)$, then
$(\rho-\sigma)\log2=2\pi i n$ and
$(\rho-\sigma)\log3=2\pi i m$ for integers $n,m$.
When $\rho\ne\sigma$, both are nonzero and elimination gives
$m\log2=n\log3$. Exponentiation gives $2^m=3^n$;
clearing negative exponents contradicts unique prime factorization.
Thus the two original prime values separate distinct roots.
Since $\lambda_p=a_p+ib_p$, the four values
$a_2,b_2,a_3,b_3$ also separate distinct roots.
Consequently, for each unordered pair, the polynomial
\[
 \Delta_{\rho\sigma}(C)=
 (a_2(\rho)-a_2(\sigma))
 +C(b_2(\rho)-b_2(\sigma))
 +C^2(a_3(\rho)-a_3(\sigma))
 +C^3(b_3(\rho)-b_3(\sigma))
\]
is a nonzero complex polynomial of degree at most three.
Also
\[
 a_p'=i(\log p)b_p,\qquad b_p'=-i(\log p)a_p.
\]
Since $a_2^2+b_2^2=1$, the first two coefficients of
\[
 D_\rho(C)=a_2'(\rho)+C b_2'(\rho)
                  +C^2a_3'(\rho)+C^3b_3'(\rho)
\]
cannot both vanish. Thus $D_\rho$ is nonzero of degree at most
three. A nonzero polynomial of degree $q$ has at most $q$ roots:
successive division by $C-c$ at distinct roots proves this by
induction on $q$. Exclude roots of $D_\rho$ only when $m_\rho>1$.
The total number of excluded complex numbers,
and hence of excluded integers, is at most $B$.
The $B+1$ displayed integers therefore contain a permitted one.
The generation assertions follow from the explicit inverse proved
in Theorem~\ref{cg:inverse-thm} below.
\end{proof}

For a completely specified finite selector one may define
\begin{equation}
 \Xi(C)=\prod_{\substack{\rho\in Z\\m_\rho>1}}D_\rho(C)
                 \prod_{\rho<\sigma}\Delta_{\rho\sigma}(C),\qquad
 c=\min\operatorname*{arg\,max}_{0\le n\le B}|\Xi(n)|.
 \tag{CG5}\label{cg:selector}
\end{equation}
The maximum is positive by the proof. This is an exact finite
formula in the supplied zeros, not a replacement by rounded
ordinates or a claim to have certified unspecified numerical inputs.
Derivatives at simple factors are unused. If a nonzero derivative
at every root is additionally wanted, test instead
$0,\ldots,3\binom r2+3r$ and include every $D_\rho$ in $\Xi$.
The sharper bound CG4 suffices for every formula below. An empty
product in CG5 is one.

Every coefficient needed in the construction has the following
explicit expression in the original prime jets. Put
$\ell_p=\log p$ and $x_{p,\rho}=p^{\rho-1/2}$. Then
\begin{equation}
\begin{split}
 h_{\rho,j}:=[t_\rho^j]H_c(\rho+t_\rho)
 =\frac1{2j!}\bigl\{
 &\ell_2^j[(1-ic)x_{2,\rho}
               +(-1)^j(1+ic)x_{2,\rho}^{-1}]\\
 +&\ell_3^j[(c^2-ic^3)x_{3,\rho}
               +(-1)^j(c^2+ic^3)x_{3,\rho}^{-1}]
 \bigr\}.
\end{split}
\tag{CG6}\label{cg:coefficients}
\end{equation}
In particular $h_{\rho,0}=z_\rho$ and $h_{\rho,1}=d_\rho$.
The full reflection identity is
$h_{\iota\rho,j}=(-1)^j\overline{h_{\rho,j}}$.
At a fixed root $z_\rho$ is real and $d_\rho$ is purely imaginary;
at a nonfixed pair the two distinct nodes are conjugate and nonreal.

\subsection{The exact confluent evaluation determinant}

Order the roots as $\rho_1,\ldots,\rho_r$ and order the rows first
by root, then by $j=0,\ldots,m_\rho-1$. Columns are indexed by
$k=0,\ldots,d-1$. Define
\begin{equation}
 E_{(\rho,j),k}=[t_\rho^j]H_c(\rho+t_\rho)^k
       =\frac1{j!}\left.\frac{d^j}{ds^j}H_c(s)^k\right|_{s=\rho}.
 \tag{CG7}\label{cg:evaluation}
\end{equation}
This is the coefficient matrix of
$P\mapsto(P(H_c))_\rho$ in the original $t_\rho$ coordinates.

\begin{theorem}[The determinant, with its orientation and constants]
\label{cg:det-thm}
For the stated row and column orders,
\begin{equation}
 \det E=
 \prod_{\rho\in Z}d_\rho^{m_\rho(m_\rho-1)/2}
 \prod_{1\le a<b\le r}
          (z_{\rho_b}-z_{\rho_a})^{m_{\rho_a}m_{\rho_b}}.
 \tag{CG8}\label{cg:det}
\end{equation}
If row $(\rho,j)$ instead contains the unscaled derivative
$\left.\frac{d^j}{ds^j}H_c(s)^k\right|_{s=\rho}$, its determinant
is CG8 multiplied by
$\prod_\rho\prod_{j=0}^{m_\rho-1}j!$.
\end{theorem}

\begin{proof}
Put $k_\rho(t)=H_c(\rho+t)-z_\rho$ and introduce the matrices
\[
 V_{(\rho,l),k}=\binom kl z_\rho^{k-l},\qquad
 (B_\rho)_{j,l}=[t^j]k_\rho(t)^l
 \quad(0\le j,l<m_\rho).
\]
The first formula is understood as zero for $l>k$; for $l=k$
its value is one, including $z_\rho=0$. Expanding a polynomial
at $z_\rho$ and substituting $k_\rho(t)$ proves
$E=\operatorname{diag}(B_\rho)V$.
Each $B_\rho$ is lower triangular, with diagonal
$1,d_\rho,d_\rho^2,\ldots,d_\rho^{m_\rho-1}$, so its determinant
is the corresponding factor in CG8. When $m_\rho=1$ the block
is $(1)$ and its determinant factor is one, even if $d_\rho=0$.

To compute $\det V$ without omitting its sign, replace the monomial
column basis by the polynomials
\[
 N_{a,j}(X)=\prod_{b<a}(X-z_{\rho_b})^{m_{\rho_b}}
                       (X-z_{\rho_a})^j,
 \qquad 0\le j<m_{\rho_a},
\]
in that order. Their degrees are exactly $0,\ldots,d-1$, and they
are monic. The column change has determinant one. At every earlier
root the whole block of jets vanishes, so the evaluation matrix
is block lower triangular. In its block at $\rho_a$, the diagonal
entry for every $j$ is
$\prod_{b<a}(z_{\rho_a}-z_{\rho_b})^{m_{\rho_b}}$.
That block is itself lower triangular, since $(X-z_{\rho_a})^j$
has no Taylor coefficient of order less than $j$.
Multiplying its diagonal entries over all $a,j$ yields the second
product of CG8 with exactly the displayed orientation.
Finally changing a Taylor-coefficient row into a derivative row
multiplies it by $j!$, which proves the last assertion.
\end{proof}

There is also a fully specified real determinant if required.
Order all fixed roots first, and then consecutive reflection pairs
$(\rho,\iota\rho)$; within each complex factor use increasing $j$.
For a fixed root take real output coordinates $i^j f_{\rho,j}$,
the coefficients in $y_\rho=t_\rho/i$. For a pair take first all
$\Re f_{\rho,j}$ and then all $\Im f_{\rho,j}$, where the second
factor is determined by reflection. The real matrix $E_\mathbb{R}$ on
real polynomial coefficients then satisfies
\begin{equation}
 \det E_\mathbb{R}=
 \left(\prod_{\rho=\iota\rho}i^{m_\rho(m_\rho-1)/2}\right)
 \left(\prod_{\{\rho,\iota\rho\},\,\rho\ne\iota\rho}
       (i/2)^{m_\rho}(-1)^{m_\rho(m_\rho-1)/2}\right)\det E.
 \tag{CG9}\label{cg:real-det}
\end{equation}
This expression is real and nonzero. Indeed the fixed-factor row
change has diagonal $i^j$. At a pair the row-change matrix is
$\left(\begin{smallmatrix}I/2&D/2\\ I/(2i)&-D/(2i)\end{smallmatrix}\right)$,
where $D=\operatorname{diag}((-1)^j)$; its determinant is the
displayed pair factor. On columns $H_c^k$ the rows are real by
CG1, establishing the reality assertion as well as the formula.

\subsection{An inverse that recovers each original Taylor coefficient}

For $m=m_\rho$, let
\begin{equation}
 q_\rho(z)=\sum_{n=1}^{m-1}v_{\rho,n}z^n,
 \qquad k_\rho(q_\rho(z))\equiv z\pmod{z^m}.
 \tag{CG10}\label{cg:inverse-coordinate}
\end{equation}
For $m=1$ this is the zero polynomial in the zero-dimensional
positive-order ideal. For $m\ge2$, all its coefficients are fixed
by the finite recursion
\begin{equation}
 \begin{split}
 v_{\rho,1}&=d_\rho^{-1},\\
 v_{\rho,n}&=-d_\rho^{-1}\sum_{l=2}^{n}h_{\rho,l}
   [z^n]\left(\sum_{j=1}^{n-1}v_{\rho,j}z^j\right)^l
                  \quad(2\le n<m).
 \end{split}
 \tag{CG11}\label{cg:inverse-recursion}
\end{equation}
To prove this, in the coefficient of $z^n$ in
$k_\rho(q_\rho(z))$, the unknown $v_{\rho,n}$ occurs only in the
linear term, with coefficient $d_\rho\ne0$. CG11 sets that
coefficient to zero for $n\ge2$, and the first line sets it to
one for $n=1$. It also follows that
\begin{equation}
 q_\rho(k_\rho(t))\equiv t\pmod{t^m}.
 \tag{CG12}\label{cg:two-sided}
\end{equation}
For a direct verification, substitution $z\mapsto k_\rho(t)$
is an invertible linear map on truncated polynomials: its matrix
is triangular with diagonal $1,d_\rho,\ldots,d_\rho^{m-1}$.
CG10 says that substitution by $q_\rho$ is a one-sided inverse;
finite-dimensional invertibility gives the other side. For $m=1$
the same assertion is the identity on constants, without a derivative
condition. Thus CG12
recovers $t$, not a replacement coordinate.

Set
\begin{equation}
 W_\rho(X)=\prod_{\sigma\ne\rho}(X-z_\sigma)^{m_\sigma},
 \qquad
 C_\rho(z)=\frac1{W_\rho(z_\rho+z)}\pmod{z^{m_\rho}}.
 \tag{CG13}\label{cg:reciprocal}
\end{equation}
The constant denominator is the exact nonzero product
$W_\rho(z_\rho)=\prod_{\sigma\ne\rho}(z_\rho-z_\sigma)^{m_\sigma}$.
For an explicit finite inverse, if
$W_\rho(z_\rho+z)=\sum_{j\ge0}w_{\rho,j}z^j$ and
$C_\rho(z)=\sum_{n=0}^{m_\rho-1}c_{\rho,n}z^n$, then
\begin{equation}
 c_{\rho,0}=w_{\rho,0}^{-1},\qquad
 c_{\rho,n}=-w_{\rho,0}^{-1}
              \sum_{j=1}^{n}w_{\rho,j}c_{\rho,n-j}
 \quad(1\le n<m_\rho).
 \tag{CG14}\label{cg:reciprocal-recursion}
\end{equation}
Coefficients of $W_\rho$ beyond its degree are zero. These are
exact finite polynomial operations at the given nodes.

\begin{theorem}[Complete inverse interpolation]\label{cg:inverse-thm}
For $f=(f_\rho(t))\in A_Z$, with
$f_\rho(t)=\sum_{j=0}^{m_\rho-1}f_{\rho,j}t^j$, define
\begin{equation}
 P_f(X)=\sum_{\rho\in Z}W_\rho(X)
    \left.
      \left[f_\rho(q_\rho(z))C_\rho(z)\right]_{<m_\rho}
    \right|_{z=X-z_\rho}.
 \tag{CG15}\label{cg:inverse}
\end{equation}
Then $\deg P_f<d$ and $P_f(H_c)=f$ in every full original local
factor. It is the unique polynomial of degree less than $d$ with
this property. In particular the exact inverse matrix is
\begin{equation}
 (E^{-1})_{k,(\rho,j)}=[X^k]\left\{
 W_\rho(X)
 \left.\left[q_\rho(z)^jC_\rho(z)\right]_{<m_\rho}
                                  \right|_{z=X-z_\rho}\right\},
 \quad 0\le k<d, 0\le j<m_\rho.
 \tag{CG16}\label{cg:inverse-matrix}
\end{equation}
For $j=0$, the factor $q_\rho(z)^0$ is one, also when $m_\rho=1$.
Moreover
\begin{equation}
 Q_H(X)=\prod_{\rho\in Z}(X-z_\rho)^{m_\rho}\in\mathbb{R}[X],\qquad
 \mathbb{R}[X]/(Q_H)\xrightarrow[\ X\mapsto H_c\ ]{\ \cong\ }A_{Z,\mathbb{R}},
 \quad
 \mathbb{C}[X]/(Q_H)\cong A_Z.
 \tag{CG17}\label{cg:real-generator}
\end{equation}
For $f\in A_{Z,\mathbb{R}}$ the polynomial $P_f$ in CG15 has real
coefficients.
\end{theorem}

\begin{proof}
Each summand in CG15 has degree at most
$(d-m_\rho)+(m_\rho-1)=d-1$.
At the node $z_\rho$, all other summands vanish modulo
$(X-z_\rho)^{m_\rho}$, while the $\rho$ summand has expansion
$f_\rho(q_\rho(z))$ modulo $z^{m_\rho}$ by CG13.
Substituting $z=k_\rho(t)$ and using CG12 gives exactly
$f_\rho(t)$ modulo $t^{m_\rho}$. This proves all original jets,
including the positive orders.
If a polynomial $P$ vanishes after substituting $H_c$ at $\rho$,
the invertible local substitution proved above says that
$P(z_\rho+z)$ vanishes modulo $z^{m_\rho}$.
Thus $(X-z_\rho)^{m_\rho}$ divides $P$.
For distinct nodes these powers are relatively prime; repeated
division, or the elementary factor theorem with multiplicities,
shows that their product $Q_H$ divides $P$.
A polynomial of degree less than $d$ divisible by $Q_H$ is zero.
This proves uniqueness and the complex isomorphism in CG17.
Applying CG15 to one coordinate vector $t_\rho^j$ gives CG16.
This is the specialized Hermite interpolation construction, with
the Chinese-remainder mechanism explicitly carried out; see the human-source
\href{https://stacks.math.columbia.edu/tag/00DT}{Stacks Project,
Lemma 10.15.4, tag 00DT} for the general theorem.

Since $H_c^\dagger=H_c$, one has
$z_{\iota\rho}=\overline{z_\rho}$, with the same multiplicity.
The coefficients of $Q_H$ are therefore real. If $P_f(H_c)=f$
and $f=f^\dagger$, then
$\overline{P_f}(H_c)=f^\dagger=f$.
Both polynomials have degree less than $d$, so uniqueness gives
$\overline{P_f}=P_f$. This proves surjectivity of the real map;
the same divisibility argument proves its kernel is $(Q_H)$.
It follows as well that $1,H_c,\ldots,H_c^{d-1}$ form a real
basis of $A_{Z,\mathbb{R}}$, and the same sequence is a complex
basis of $A_Z$.
\end{proof}

The auxiliary polynomial $Q_H$ is the newly derived relation in
the variable $X$. It does not replace $g(s)$, the original packet
polynomial $h(s)=c_h\prod(s-\rho)^{m_\rho}$, its chosen leading
coefficient $c_h$, or any original local unit $g/h$.
All of those data remain unchanged when CG15 is used as the
input vector of the original primitive lift.

For further direct checks on the real descent, let
$P_{\rho,j}$ denote the polynomial in braces in CG16.
Uniqueness and CG1 give
\begin{equation}
 P_{\iota\rho,j}(X)=(-1)^j\overline{P_{\rho,j}(X)}.
 \tag{CG18}\label{cg:paired-inverse}
\end{equation}
Indeed dagger sends the original jet vector $t_\rho^j$ supported
at $\rho$ to $(-1)^j t_{\iota\rho}^j$ supported at $\iota\rho$.
This proves CG18 and verifies every reflection sign in the
inverse formula.

\subsection{Finite Laurent words and every original prime factor}

In $\mathbb{C}[X_2^{\pm1},X_3^{\pm1}]$ define
\begin{equation}
 \mathcal H_c=A_+X_2+A_-X_2^{-1}+B_+X_3+B_-X_3^{-1},
 \quad A_\pm=\frac{1\mp ic}{2},\qquad
 B_\pm=\frac{c^2\mp ic^3}{2}.
 \tag{CG19}\label{cg:laurent-h}
\end{equation}
Under the original two-prime map $X_p\mapsto\lambda_p$, this is
exactly $H_c$. For $k\ge0$ its complete expansion is
\begin{equation}
 \mathcal H_c^k=
 \sum_{\substack{u_+,u_-,v_+,v_-\ge0\\u_++u_-+v_++v_-=k}}
 \frac{k! A_+^{u_+}A_-^{u_-}B_+^{v_+}B_-^{v_-}}
      {u_+!u_-!v_+!v_-!}
 X_2^{u_+-u_-}X_3^{v_+-v_-}.
 \tag{CG20}\label{cg:multinomial}
\end{equation}
This follows by choosing one of the four displayed terms in
each of the $k$ factors; the multinomial coefficient counts
exactly the choices with those four counts.

\begin{proposition}[A finite bound retaining all jets and original scales]
\label{cg:words-thm}
Every $f\in A_Z$ has the Laurent representative
$P_f(\mathcal H_c)$ supported on
\begin{equation}
 \{(u,v)\in\mathbb{Z}^2:|u|+|v|\le d-1\}.
 \tag{CG21}\label{cg:diamond}
\end{equation}
There are exactly $2d^2-2d+1$ possible distinct Laurent words
in this set. If $f\in A_{Z,\mathbb{R}}$ then $P_f$ is real and the
Laurent coefficient at $(-u,-v)$ is the complex conjugate of
the coefficient at $(u,v)$.
For every integer pair, its action on the original packet is
\begin{equation}
 X_2^uX_3^v\longmapsto
 M_{2^{u(s-1/2)}3^{v(s-1/2)}}
 =2^{-u/2}3^{-v/2}\Pi_2^u\Pi_3^v,
 \qquad u,v\in\mathbb{Z}.
 \tag{CG22}\label{cg:original-words}
\end{equation}
For any further rational prime $p$, define the entire functions
$a_p,b_p$ by CG2. Their full jets have unique real polynomials
$P_{a_p},P_{b_p}$ given by CG15, each of degree at most $d-1$.
Then its original action is exactly
\begin{equation}
 \Pi_p=p^{1/2}M_{\,P_{a_p}(H_c)+iP_{b_p}(H_c)}.
 \tag{CG23}\label{cg:all-primes}
\end{equation}
In particular the same finite Laurent support bound suffices
for every other original prime action on the given full packet.
\end{proposition}

\begin{proof}
Each monomial in CG20 has
$|u_+-u_-|+|v_+-v_-|\le k$.
Since $P_f$ has degree at most $d-1$, its expansion is supported
in CG21. For $D\ge0$, the shell $|u|+|v|=j$ has $4j$ points
for $j\ge1$, while the central shell has one. The total is
$1+\sum_{j=1}^D4j=1+2D(D+1)$; setting $D=d-1$ gives the
claimed exact cardinality. The Laurent involution conjugates
coefficients and sends $X_p$ to $X_p^{-1}$.
In CG19, $A_- =\overline{A_+}$ and
$B_-=\overline{B_+}$, so $\mathcal H_c$ is fixed.
A real polynomial in it is also fixed, proving the coefficient
relation.

The original operators $\Pi_p$ are invertible on $A_Z$ because
the constant term $p^\rho$ is nonzero at every root.
Thus $X_p=p^{-1/2}\Pi_p$ and
$X_p^{-1}=p^{1/2}\Pi_p^{-1}$ hold with their full jets.
Taking integer powers proves CG22, including negative exponents.
Finally $a_p+ib_p=\lambda_p$ as an entire identity, so CG15
applied to its two real components recovers the full Taylor
series of $\lambda_p$ at every root. Multiplication by the
retained original factor $p^{1/2}$ proves CG23.
\end{proof}

One may equivalently apply CG15 directly to
\begin{equation}
 (\lambda_p)_\rho(t)=p^{\rho-1/2}
              \sum_{j=0}^{m_\rho-1}\frac{(\log p)^j}{j!}t^j,
 \qquad
 (p^s)_\rho(t)=p^{\rho}
              \sum_{j=0}^{m_\rho-1}\frac{(\log p)^j}{j!}t^j.
 \tag{CG24}\label{cg:other-prime-jets}
\end{equation}
Linearity of CG15 then gives
$P_{p^s}=p^{1/2}P_{\lambda_p}$ and
$P_{\lambda_p}=P_{a_p}+iP_{b_p}$.
No prime eigenvalue has been replaced by its modulus.

\subsection{Edge cases and the exact scope of the bounds}

The case $r=1,m_\rho=1$ gives $d=1$, $E=(1)$,
$W_\rho=1$, $q_\rho=0$, and $P_f=f_\rho$.
The empty determinant products in CG8 equal one.
CG21 consists only of $(0,0)$ and has cardinality one.
The selector CG4 gives $B=0$ and permits $c=0$ without a
derivative test.
Since a one-root reflection-stable packet is fixed by reflection,
its real algebra in this case is $\mathbb{R}$.

For $r=1$ with any larger multiplicity, CG4 gives $B=3$.
Distinct-node conditions
are empty and a nonzero derivative still gives the full truncated
factor. The determinant is exactly
$d_\rho^{m_\rho(m_\rho-1)/2}$. The local recursion retains each
of its $m_\rho-1$ positive-order coefficients.
No argument assumes the multiplicities are equal at different
orbits; equality is needed only at the two members of one
reflection orbit, as provided by the original reflection.

A chosen integer may be zero. In that case $B_\pm=0$, and terms
with positive corresponding multinomial counts vanish. This can
only reduce Laurent support; it does not invalidate the bound.
Some $z_\rho$ may be zero. The determinant and inverse use only
differences of distinct nodes and nonzero $d_\rho$ at nonsimple
roots, so this
causes no division by zero. Interpret $z^0=1$ and constant
polynomial coefficients directly, including the $m_\rho=1$
case in CG16.

The degree bound $d-1$ is optimal for representing every element
by polynomials in this one generator: polynomials of degree at
most $D<d-1$ have dimension $D+1<d$, so cannot span $A_Z$.
No optimality assertion is made for arbitrary two-variable
Laurent representatives. The $2d^2-2d+1$ words are a spanning
list with many relations, not a basis of a $d$-dimensional
space. Products of two unreduced representatives may have
Laurent degree $2d-2$, but their packet product has a new
representative of degree at most $d-1$ by CG15 or reduction
modulo $Q_H$; no original jet is lost in this reduction.

The candidate count in CG4 depends on the number of distinct
roots, whereas the interpolation and Laurent-degree counts
depend on the sum of full multiplicities. Replacing $d$ by $r$
in those latter bounds would discard the nilpotent directions
when a multiplicity exceeds one.
The finite selector produces a generator for each supplied
finite packet. It does not assert the existence of a single
integer that works simultaneously for all finite packets, nor
does it assert a bound independent of the packet size.



\section{The actual four dilation terms and their finite theta relation}
For every positive real $q$ set
\begin{equation}
 U_qF(x)=F(x/q),\qquad V_q=q^{-1/2}U_q,\qquad
 \MM V_qF(s)=q^{s-1/2}\MM F(s).\label{eq:Vq}
\end{equation}
The same formulas act on the Schwartz source and commute with $\Theta$.
Substitution in $dx$ gives $\|U_qF\|^2=q\|F\|^2$, so $V_q$ is unitary
and $V_q^*=V_{q^{-1}}$. The original prime action is exactly
$U_p=p^{1/2}V_p$, not $V_p$ alone.
Use the real integer $c$ constructed in Section 2 and the four coefficients
\begin{equation}
 \alpha_2=\frac{1-ic}{2},\quad\alpha_{1/2}=\frac{1+ic}{2},\quad
 \alpha_3=\frac{c^2-ic^3}{2},\quad\alpha_{1/3}=\frac{c^2+ic^3}{2}.
 \label{eq:alphas}
\end{equation}
Their original operator and its finite-packet image are
\begin{equation}
 \HS=\sum_{q\in\{2,1/2,3,1/3\}}\alpha_q V_q,
 \qquad M=M_{H_c}\in\operatorname{End}(A_h).
 \label{eq:operator}
\end{equation}
Each pair of coefficients is conjugate, proving self-adjointness on
$L^2(dx)$. On the actual Mellin line its multiplier is the real function
\begin{align}
 x_c(t):=H_c(1/2+it)
 &=\cos(t\log2)+c\sin(t\log2)\nonumber\\
 &\quad+c^2\cos(t\log3)+c^3\sin(t\log3),\nonumber\\
 |x_c(t)|&\le K_c:=(1+c^2)^{3/2},\qquad\|\HS\|\le K_c.
 \label{eq:bound}
\end{align}
Each sine--cosine pair has amplitude respectively
$\sqrt{1+c^2}$ and $c^2\sqrt{1+c^2}$, proving this bound.
We need no assertion of equality in the operator norm.

\subsection{The coefficient of the discrepancy, including all zero derivatives}
Let $M_s$ multiply by $s$ in $A_h$. Exact division gives the scalar
\begin{equation}
 \kappa_h(u)=\frac{sR_u-R_{M_su}}{h}
            =\frac{[s^{d-1}]R_u}{c_h},\qquad
 D\PP(u)-\PP(M_su)=\Theta(\kappa_h(u)a(D)\phi_*).
 \label{eq:kappa}
\end{equation}
The numerator vanishes modulo $h$ and has degree at most $d$, giving
the first two equalities. Taking Mellin transforms gives the last one.
There is a useful direct expression in the original local units:
\begin{equation}
 \kappa_h(u)=\sum_{\rho\in Z}\Res_\rho
                     \frac{u(s)}{a(s)g(s)}\dd s.\label{eq:kappa-residue}
\end{equation}
Represent $u$ by its local jets. The residues of $R_u/h$ and
$u/(ag)$ coincide at each selected zero, since $[aU_h]R_u=u$.
The sum of the residues of $R_u/h$ is its coefficient of $s^{-1}$
at infinity, namely $[s^{d-1}]R_u/c_h$.
This proves \eqref{eq:kappa-residue}, and records exactly where the
unchanged $c_h$ cancels against the unit in $R_u$.
If $g(\rho+t)=t^{m_\rho}\gamma_\rho(t)$, define
\begin{equation}
 B_{\rho,j}(u)=[t^j]\frac{u_\rho(t)}{a(\rho+t)\gamma_\rho(t)}.
\end{equation}
Expanding the exponential in the residue then yields the finite formula
\begin{equation}
 \kappa_h(e^{vM_s}u)=\sum_\rho e^{v\rho}
       \sum_{k=0}^{m_\rho-1}\frac{v^k}{k!}B_{\rho,m_\rho-1-k}(u).
 \label{eq:kappa-exp}
\end{equation}
No derivative of the actual zero unit has been replaced by its value.

\begin{theorem}[An exact source relation for the four original dilations]
For $L\in\R$ use an oriented integral and set
\begin{align}
 \psi_L(u)&=\int_0^L\kappa_h(e^{vM_s}u)
           U_{e^{L-v}}(a(D)\phi_*)\dd v,\nonumber\\
 \Psi_c(u)&=\sum_q\alpha_q q^{-1/2}\psi_{\log q}(u).
 \label{eq:psi}
\end{align}
Then $\Psi_c(u)\in V$, its theta transform has endpoint Mellin values
zero, and
\begin{equation}
 \HS\PP(u)-\PP(Mu)=\Theta\Psi_c(u).\label{eq:homotopy}
\end{equation}
The interval for $q<1$ retains its negative orientation.
\end{theorem}
\begin{proof}
Differentiate $U_{e^{L-v}}\PP(e^{vM_s}u)$ in $v$. Equation
\eqref{eq:kappa} gives derivative
\[
 -\Theta[\kappa_h(e^{vM_s}u)U_{e^{L-v}}a(D)\phi_*].
\]
Integrating from $0$ to $L$, with either sign of $L$, proves
$U_{e^L}\PP(u)-\PP(e^{LM_s}u)=\Theta\psi_L(u)$.
Multiply by $q^{-1/2}$ and the four coefficients to get
\eqref{eq:homotopy}. The integrands are smooth Schwartz families on
compact intervals. Their source conditions are preserved by positive
dilation. Their theta transforms are dilates of $ag$, which vanish
at $s=0,1$. This proves every asserted space membership.
\end{proof}

This is also a quantitative relation, in the same inherited observation.
Set
\[
 C_*:=\|\Theta(a(D)\phi_*)\|_{L^2(dx)}<\infty.
\]
Then
\begin{equation}
 \|\Theta\Psi_c(u)\|\le C_*\sum_q|\alpha_q|
   \int_{\min(0,\log q)}^{\max(0,\log q)}
      e^{-v/2}|\kappa_h(e^{vM_s}u)|\dd v.\label{eq:cost}
\end{equation}
Indeed the norm of the dilated seed is $e^{(L-v)/2}C_*$;
the factor $q^{-1/2}=e^{-L/2}$ cancels precisely its $e^{L/2}$.
The triangle inequality for the finite sum and each integral proves
\eqref{eq:cost}. Its integrand is the finite exponential polynomial
in \eqref{eq:kappa-exp}, made from actual zero derivatives.

\section{Finite prime-word representatives, with exact source and Weil metrics}
Put $F_*:=\PP(1)$ and define
\begin{equation}
 F_j=\HS^jF_*,\qquad\Phi(P)=P(\HS)F_*,\qquad
 \widehat F_*(s):=\MM F_*(s)=a(s)U_h(s)R_1(s).
 \label{eq:word-family}
\end{equation}
These functions are in $\B$, have both endpoint moments zero, and
have full selected jets $H_c^j$. Their transforms vanish at every
other actual zero to its full order. Powers through $j=d-1$ give a
basis of the entire arithmetic packet by Section 2.
For every $j\ge1$ the exact comparison with the original section is
\begin{equation}
 \HS^j\PP(u)-\PP(M^ju)=\Theta\mathcal D_j(u),\quad
 \mathcal D_j(u)=\sum_{k=0}^{j-1}\HS^{j-1-k}\Psi_c(M^ku),
 \quad\mathcal D_0=0.\label{eq:powers}
\end{equation}
The source action of $\HS$ in this formula uses the same four
dilations. Induction using commutation with $\Theta$ proves
\eqref{eq:powers}, including its order and signs. Extend
$\mathcal D_P=\sum_jp_j\mathcal D_j$ linearly in $P(X)=\sum_jp_jX^j$.

Let
\begin{equation}
 \nu(X)=\prod_\rho(X-z_\rho)^{m_\rho}\in\R[X],
 \qquad z_\rho=H_c(\rho).
 \label{eq:nu}
\end{equation}
This is the minimal polynomial of $M$, not a polynomial obtained by
removing multiplicities. The exact relation calculation is
\begin{equation}
 [\Phi(P)]=0\text{ in }Q
 \quad\Longleftrightarrow\quad\nu\mid P.
 \label{eq:relation-ideal}
\end{equation}
Necessity follows from the full jet isomorphism. For sufficiency,
$P(M)=0$ and \eqref{eq:powers} give the actual source primitive
$\Phi(P)=\Theta\mathcal D_P(1)$.
At a fixed original nonempty support label, this polynomial quotient
sends killed coefficients to that label's supported zero. Every
support label is retained. Applying $G_L$ to the unital algebra
isomorphism of Section 2 therefore gives the same label-preserving
semiring morphism; no support stratum is removed.

Define the finite positive measure $\mu$ by its integrals:
\begin{equation}
 \int f(x)\dd\mu(x)=\frac1{2\pi}\int_\R
  f(x_c(t))|a(1/2+it)U_h(1/2+it)R_1(1/2+it)|^2\dd t.
 \label{eq:measure}
\end{equation}
The real spectral integration variable in $\mu$ is not the positive
source coordinate inside $F(x)$; their connecting map is exactly
$t\mapsto x_c(t)$ after the Mellin transform.
Its support lies in $[-K_c,K_c]$, and its total mass is
$\|F_*\|^2$, not one. In particular
\begin{align}
 \langle\Phi(P),\Phi(Q)\rangle_{L^2(dx)}
     &=\int\overline{P(x)}Q(x)\dd\mu(x),\nonumber\\
 \mu_j&=\int x^j\dd\mu(x),\qquad
 G^{\rm word}_{jk}=\mu_{j+k}.\label{eq:source-gram}
\end{align}
Equation \eqref{eq:plancherel} and the real multiplier
\eqref{eq:bound} prove this identity. The source Gram is positive
definite at every finite degree: $\widehat F_*$ is not the zero entire
function, so its real-line zeros are isolated; $x_c$ is a nonconstant
real analytic function. The latter follows from independence of its
four distinct exponential frequencies, since $\alpha_2\ne0$.
On an interval where $x_c'\ne0$ and $\widehat F_*\ne0$, the measure
has positive mass in every subinterval of its image. A nonzero
polynomial cannot vanish on that interval. This proves the asserted
strict positivity and also that the map $\Phi$ before quotient is
injective in every polynomial degree.

The whole Weil form, on these same representatives, is a different
explicit matrix:
\begin{equation}
 T^{\rm word}_{jk}=\mathcal W_g(F_j,F_k)
       =\sum_{\rho\in Z}m_\rho z_\rho^{j+k}.
 \label{eq:Weil-gram}
\end{equation}
Indeed all omitted-zero summands are identically zero by
\eqref{eq:word-family}; and
$\overline{H_c(\iota\rho)}=H_c(\rho)$ gives the displayed powers.
Thus \eqref{eq:source-gram} and \eqref{eq:Weil-gram} calculate two
forms on the identical functions, with an explicit identity map
between their underlying vectors. The finite trace has radical the
positive-order jet ideal. Its matrix on values is positive $m_\rho$
at a reflection-fixed root and
$\left(\begin{smallmatrix}0&m_\rho\\m_\rho&0\end{smallmatrix}\right)$
at a two-element reflection orbit. These statements follow by writing
the two terms of the reflected sum. The raw pairing
\[
 R_h(u,v)=\sum_\rho\Res_\rho\frac{u^\dagger v}{g}\dd s
\]
still detects all jets: its local coefficient matrix has nonzero
antidiagonal $(-1)^j/\gamma_\rho(0)$. It is perfect, and insertion
of $g'$ gives the trace because $g'/g=m_\rho/t+\gamma_\rho'/\gamma_\rho$.
This proves the connecting map between the raw and contracted forms.

\section{Bounded moment dynamics and the attained quotient volume}
For $\theta\in\C$ set
\begin{equation}
 \begin{gathered}
 Z(\theta)=\int e^{\theta x}\dd\mu(x),\qquad
 Z_\nu(\theta)=\int |\nu(x)|^2e^{\theta x}\dd\mu(x),\\
 D_n=\det[Z^{(i+j)}]_{i,j<n},\qquad
 B_n=\det[Z_\nu^{(i+j)}]_{i,j<n}.
 \end{gathered}\label{eq:determinants}
\end{equation}
Empty determinants are $D_0=B_0=1$. Both transforms are entire, since
on every compact set of $\theta$ the integrand and all its derivatives
are bounded by a constant times the finite measure. Their derivatives
are exactly the indicated moments. For real $\theta$, every
positive-order determinant is positive; multiplication by $\nu$ loses
only finitely many real points from the interval used to prove strict
positivity. In particular no atom assumption is needed for this step.

For $n\ge d-1$, the original word-degree cutoff has the exact sequence
\begin{equation}
 0\longrightarrow\C[X]_{\le n-d}
 \xrightarrow{\times\nu}\C[X]_{\le n}
 \longrightarrow\C[X]/(\nu)\longrightarrow0.\label{eq:sequence}
\end{equation}
Set $\C[X]_{\le-1}=0$. Give the middle space precisely the measure
$e^{\theta x}\mu$ and the quotient its unique least-norm
representative metric. At $\theta=0$ this is the inherited source
metric of the displayed prime words. In the quotient basis
$1,X,\ldots,X^{d-1}$, write this Gram matrix as $G_n(\theta)$.
Then
\begin{equation}
 \mathcal V_n(\theta):=\det G_n(\theta)
       =\frac{D_{n+1}(\theta)}{B_{n-d+1}(\theta)}.
 \label{eq:volume}
\end{equation}
To prove it, use the middle-space basis
$1,\ldots,X^{d-1},\nu,X\nu,\ldots,X^{n-d}\nu$.
The monic polynomial $\nu$ gives change-of-basis determinant one.
The relation block has Gram $B_{n-d+1}$. Orthogonally subtracting
that block gives the Schur complement, exactly the least-norm
quotient Gram. Taking determinants proves \eqref{eq:volume}.
The auxiliary observation has an exact isometry to the original
word measure: $f(x)\mapsto e^{\theta x/2}f(x)$ for real $\theta$.
On the source Hilbert space this is $e^{\theta\HS/2}$, with both it
and its inverse bounded by $e^{|\theta|K_c/2}$. Its image of the
polynomial subspace is specified by this formula; it is not assumed
to preserve a polynomial cutoff.
These are also actual operators on $\B$: for its defining seminorm
$p_{b,j}$ one has $p_{b,j}(V_qF)=q^{b-1/2}p_{b,j}(F)$, hence
$p_{b,j}(\HS F)\le C_b p_{b,j}(F)$ with
$C_b=\sum_q|\alpha_q|q^{b-1/2}$.
The exponential series converges in every such seminorm. The
same dilation estimate on the Schwartz source proves preservation
of its source conditions and theta relations. Taking Mellin
transforms termwise gives multiplication by $e^{\theta H_c(s)/2}$.
For clarity, its full packet action is $E_\theta=e^{\theta M/2}$:
the transported seed has jets $e^{\theta H_c/2}$, not $1$.
If these transported representatives are then re-expressed in the
same original jet coordinates, their quotient-volume determinant is
\begin{equation}
 \widetilde{\mathcal V}_n(\theta)
 =e^{-\theta\sum_\rho m_\rho z_\rho}\mathcal V_n(\theta).
 \label{eq:tilt-coordinate}
\end{equation}
Indeed their Gram in fixed jets is
$E_{-\theta}^*G_nE_{-\theta}$; its determinant changes by
$|\det E_{-\theta}|^2=e^{-\theta\Tr M}$.
The trace is real by reflection. This retained linear logarithmic
term has zero second derivative, so the curvature below is unchanged.

\begin{theorem}[Cutoff-independent arithmetic recurrence bounds]
For each of the two positive measures $e^{\theta x}\mu$ and
$e^{\theta x}|\nu(x)|^2\mu$, let $p_j$ be its monic orthogonal
polynomial and $h_j$ its squared norm. Its three-term recurrence is
\begin{equation}
 xp_j=p_{j+1}+b_jp_j+a_jp_{j-1},\qquad
 a_j=h_j/h_{j-1},\quad a_0=0.
 \label{eq:recurrence}
\end{equation}
For every real $\theta$ and $j\ge1$,
\begin{equation}
 0<a_j\le K_c^2=(1+c^2)^3.\label{eq:uniform}
\end{equation}
Writing $a_j^{\rm src}$ and $a_j^{\rm rel}$ for the two sequences,
the exact volume identity is
\begin{equation}
 (\log\mathcal V_n)''=a_{n+1}^{\rm src}-a_{n-d+1}^{\rm rel},
 \qquad |(\log\mathcal V_n)''|\le K_c^2.
 \label{eq:curvature}
\end{equation}
At $n=d-1$ the second term is zero, without a negative-index determinant.
\end{theorem}
\begin{proof}
Orthogonality and symmetry of multiplication by $x$ show that in
$xp_j$ only degrees $j+1,j,j-1$ occur; pairing with $p_{j-1}$ gives
the displayed $a_j$. The coefficient of $p_j/\sqrt{h_j}$ in
$x p_{j-1}/\sqrt{h_{j-1}}$ is $\sqrt{h_j/h_{j-1}}$.
Multiplication by $x$ has norm at most $K_c$ on each of the two
measures. Cauchy--Schwarz proves \eqref{eq:uniform}.

Differentiating orthogonality with respect to the real parameter,
including the derivative $x e^{\theta x}$ of the measure, gives
$\partial_\theta p_j=-a_jp_{j-1}$ and
$\partial_\theta\log h_j=b_j$.
For precision, pairing $\partial_\theta p_j$ with a degree $k<j-1$
gives zero because $xp_k$ has degree below $j$; pairing with
$p_{j-1}$ gives $-h_j$, yielding exactly $-a_j$.
The derivative terms from $p_j$ in $h_j'$ vanish by orthogonality.
Now differentiate $b_j=\langle xp_j,p_j\rangle/h_j$.
The measure derivative contributes
$a_{j+1}+b_j^2+a_j$, the two polynomial derivatives contribute
$-2a_j$, and the denominator contributes $-b_j^2$.
Thus $b_j'=a_{j+1}-a_j$. Since $D_n=\prod_{j<n}h_j$,
telescoping gives
\[
 (\log D_n)''=a_n=\frac{D_{n+1}D_{n-1}}{D_n^2}
 \quad(n\ge1).
\]
The same proof gives the identity for $B_n$. Subtracting their
logarithms in \eqref{eq:volume} proves \eqref{eq:curvature}.
Both recurrence terms lie in $[0,K_c^2]$, so their difference has
absolute value at most $K_c^2$, not twice that bound.
\end{proof}

These are the classical Hankel--Toda identities, here derived for the
explicit measure \eqref{eq:measure}; their general orthogonal-polynomial
provenance is recorded in \cite{Toda}. This is a concrete continuation of
the programme's source/relation-volume calculation, not a new choice of
positive metric on $A_h$. The earlier Euler-polynomial family uses
$P(D)F_h$ and an unbounded Mellin coordinate. Here the family is
$P(\HS)\PP(1)$, a finite sum of dilations. Equation \eqref{eq:powers}
is their exact cohomological comparison with the same original section;
\eqref{eq:source-gram} supplies the new family's actual metric.
No bound proved for this family is silently transferred to the earlier
Euler cutoffs. For the identity map on the finite packet, its two
representative maps differ by the displayed theta source relation.

\section{The relation cost retains every spectral discrepancy}
For a polynomial $P$, \eqref{eq:powers} gives
\begin{equation}
 P(\HS)\PP(u)=\PP(P(M)u)+\Theta\mathcal D_P(u).
 \label{eq:polynomial-defect}
\end{equation}
Let $z\in\C$ and $m\ge1$, and put
\begin{equation}
 d_K(z)=\sqrt{(\Im z)^2+\max\{|\Re z|-K_c,0\}^2}.
\end{equation}
By multiplication on the Mellin line, for every $u\in A_h$,
\begin{equation}
 \|\Theta\mathcal D_{(X-z)^m}(u)\|
 \ge d_K(z)^m\|\PP(u)\|-\|\PP((M-z)^mu)\|.
 \label{eq:residual}
\end{equation}
Indeed $|x_c(t)-z|\ge d_K(z)$ pointwise, so the left-hand side of
\eqref{eq:polynomial-defect} has norm at least
\[
 d_K(z)^m\|\PP(u)\|.
\]
The triangle inequality proves
\eqref{eq:residual}. On the actual generalized eigenspace
$\ker(M-z)^m$, the last term is exactly zero. This calculates a
necessary nonzero theta-relation norm whenever that eigenspace has
$d_K(z)>0$. It does not identify an unbounded algebraic quotient
with a bounded Hilbert quotient.

There is a full matrix equality behind this bound. In any fixed
coordinate basis of $A_h$, write $S$ for the column map $\PP$,
$G=S^*S$ for its inherited positive Gram, and
$E=\HS S-SM=\Theta\Psi_c$ for the actual discrepancy columns.
Self-adjointness of $\HS$ gives exactly
\begin{equation}
 S^*E-E^*S=M^*G-GM.\label{eq:matrix-defect}
\end{equation}
Expanding the two left terms proves the formula; the common term
$S^*\HS S$ cancels. Every coefficient is either an original source
inner product or the explicit finite operator. This is the precise
map relating a self-adjoint source operator to its possibly
non-self-adjoint algebraic finite action.

\section{An original-zeta numerical example and exact finite checks}
Take the actual first two positive critical ordinates $\gamma_1,\gamma_2$,
their negative partners, and $c_h=1$. With
\[
 h(s)=\bigl((s-1/2)^2+\gamma_1^2\bigr)
      \bigl((s-1/2)^2+\gamma_2^2\bigr),\qquad c=1,
\]
the numerical generator values, in the order
$-\gamma_1,+\gamma_1,-\gamma_2,+\gamma_2$, are
\[
 (-1.729633532453,\ -2.100994546437,\
   -0.885009068584,\ -0.856555247283).
\]
These decimal observations are not used to certify the general theorem
or the simplicity of a zero. The original low-height human source is
Platt--Trudgian \cite{Platt}; the numerical computation here uses
\texttt{mpmath.zetazero}, retaining its status as numerical evidence.
For these simple nodes the matrix has rows $(1,z,z^2,z^3)$ and
numerical determinant $-0.0117911416732783$.

The exact unit-inverse polynomial is obtained without discarding its
large coefficients. Put $\rho_j=1/2+i\gamma_j$ and
\[
 b_j=\frac{h'(\rho_j)}{\rho_j(\rho_j-1)g'(\rho_j)},\qquad
 r_2=\frac{b_1-b_2}{\gamma_2^2-\gamma_1^2},\qquad
 r_0=b_1+r_2\gamma_1^2.
\]
Then $R_1(s)=r_0+r_2(s-1/2)^2$: it has exactly the four required
values $[aU_h]^{-1}$, by conjugation and reflection, and degree below
four. Numerically $r_0\simeq512549.3874122573$ and
$r_2\simeq2627.36745870685$. The resulting source Gram is
\[
 \begin{pmatrix}
 16369.56790&201.64114&11653.99632&-6605.01598\\
 201.64114&11653.99632&-6605.01598&27600.94756\\
 11653.99632&-6605.01598&27600.94756&-29905.12371\\
 -6605.01598&27600.94756&-29905.12371&96519.71793
 \end{pmatrix}.
\]
Its first three monic recurrence coefficients are numerically
\[
 (a_1,a_2,a_3)\simeq(0.7117788668,\ 1.3213160586,\ 1.4695672533),
 \qquad a_j\le8\text{ for every degree by \eqref{eq:uniform}}.
\]
In particular the source unit has squared norm about $16369.56790$,
whereas its full Weil value is exactly $4$ for this simple four-root
packet. This is a calculation on the same original representative,
not a rescaling to hide the difference.

\texttt{calculate\_example.py} evaluates the source at 40 decimal
working digits, then uses double-precision composite Gauss--Legendre
quadrature on $0\le t\le80$, exploiting the proved evenness of this
example's density. Orders 12 and 20 on intervals of length two give
maximum relative moment discrepancy $3.95\times10^{-14}$ using the
comparison denominator $\max(1,|\mu_j|)$. This agreement is not an
interval certificate or a bound for the omitted tail. All numerical
matrices and coefficients are in \texttt{NUMERICAL\_EXAMPLE.json}.
\texttt{check\_identities.py} separately verifies exact rational and
symbolic finite identities, including repeated jets, nonunit $c_h$,
the oriented dilation formula, and the endpoint of the determinant
ratio. It does not certify the transcendental numerical example.

\begin{figure}[p]
 \centering\includegraphics[width=\textwidth]{two_prime_example.pdf}
 \caption{The original two-prime symbol, primitive unit density,
 distinct packet values, and inherited source Gram. These are
 numerical observations, not certified bounds. The exact operator
 and constants are \eqref{eq:alphas}--\eqref{eq:bound}; the actual
 density and full matrix are \eqref{eq:measure}--\eqref{eq:source-gram}.
 The full-jet mechanism is proved in Section 2, while the bounded
 recurrence and determinant dynamics are proved in Section 5.
 Human source locators: Meyer \cite{Meyer}, Theorem
 \texttt{the:Lap\_range}; Stacks \cite{Stacks}, tag 00DT.
 Reproducible figure source: \texttt{calculate\_example.py}.}
\end{figure}

\section*{What this calculation supplies}
There is now a finite construction of every full packet class using
the actual primes 2 and 3, an exact original theta primitive for
every discrepancy, and an inherited arithmetic word metric with
cutoff-independent recurrence and curvature bounds. The construction
retains the original $e$-relations through \eqref{eq:relation-ideal},
not through a numerical replacement of the supported zeros.
This supplies directly evaluable inputs for the programme's
source/relation determinant method. It neither proves the Riemann
hypothesis nor provides a counterexample. The new observation is
joined to the old section by a proved cochain formula, not declared
equal to a previously chosen metric or cutoff.

\begin{thebibliography}{9}
\bibitem{Meyer} Ralf Meyer, \emph{A spectral interpretation for the zeros
of the Riemann zeta function},
\href{https://arxiv.org/abs/math/0412277}{arXiv:math/0412277}.
Original author source, Theorem \texttt{the:Lap\_range} and
Corollary \texttt{cor:Dminus\_eigenvalues}; read lines 641--740
of the indexed local source. This is the human context for the
original Mellin quotient and multiplicities.
\bibitem{Stacks} The Stacks Project Authors,
\href{https://stacks.math.columbia.edu/tag/00DT}{Lemma 10.15.4,
Chinese remainder, tag 00DT}. Original \texttt{algebra.tex},
\texttt{lemma-chinese-remainder}, complete proof read.
The determinant and inverse formulas here are specialised
Hermite-interpolation calculations proved in Section 2.
\bibitem{Toda} A. I. Aptekarev, A. Branquinho and F. Marcell\'an,
\emph{Toda-type differential equations
for the recurrence coefficients of orthogonal polynomials and Freud
transformation}, Journal of Computational and Applied Mathematics
 78 (1997), 139--160, \href{https://doi.org/10.1016/S0377-0427(96)00138-0}
 {doi:10.1016/S0377-0427(96)00138-0}.
Historical provenance for the orthogonality/Toda mechanism; the
identities used in this note are proved directly in Section 5.
\bibitem{Platt} David Platt and Timothy Trudgian,
\emph{The Riemann hypothesis is true up to $3\cdot10^{12}$},
\href{https://arxiv.org/abs/2004.09765}{arXiv:2004.09765}.
Human provenance for the established low-height zero computation,
not a certificate for the numerical quadrature performed here.
\end{thebibliography}
\end{document}
